\BourbakiExercisesSection

\begin{enumerate}
\def\labelenumi{\arabic{enumi}.}
\item
  枚举长度 \(\leqslant 4\) 的、位于 \(\mathbf{M}(\mathbf{X})\) 中的元素（此时 \(\mathbf{X}\) 由单个元素组成）。
\item
  设 X 为一个集合，M 为广群 M(X)、Mo(X) 或 \(\mathrm{N}^{(X)}\) 之一。证明 M 的每个自同构都使 X 稳定。由此推出一个从 \(\mathfrak{S}_{X}\) 到群 Aut(M) 上的同构。
\end{enumerate}

证明 M 的自同态与从 X 到 M 的映射一一对应。

\begin{enumerate}
\def\labelenumi{\arabic{enumi}.}
\setcounter{enumi}{2}
\tightlist
\item
  设 X 为一个集合， \(M_{a}\) 为在单元素集 \(\{a\}\) 上构造的自由广群。设 \(\rho\) 表示将 X 的每个元素都映到 a 的、从 \(\mathrm{M}(\mathrm{X})\) 到 \(M_{a}\) 中的同态；另一方面，设 \(\lambda: \mathrm{M}(\mathrm{X}) \to \mathrm{Mo}(\mathrm{X})\) 为第 9 号中定义的同态。证明映射 \(w \mapsto (\lambda(w), \rho(w))\) 是从 \(\mathrm{M}(\mathrm{X})\) 到 \(\mathrm{Mo}(\mathrm{X}) \times \mathrm{M}_{a}\) 的由 u 与 v 具有相同长度的有序对 \((u, v)\) 所组成的子广群上的同构。
\end{enumerate}

¶ 4. 设 X 为一个集合，c 为不属于 X 的元素。记 Y = X ∪ \{c\}。设 M 为以 Mo(Y) 为底层集合、合成律由 (u, v) ↦ cuv 给出的广群。设 L 为由 X 生成的 M 的子广群，并设 ε 为从 Y 到 Z 的映射，它在 X 上取值 -1，在 c 处取值 1。

\begin{enumerate}
\def\labelenumi{(\alph{enumi})}
\tightlist
\item
  证明 L 是由满足下列两个条件的字 \(w = y_{1} \ldots y_{p}\) , \(y_{i} \in Y\) 组成的 M 的子集：
\end{enumerate}

\[
(i) \varepsilon (y _ {1}) + \dots + \varepsilon (y _ {p}) = - 1
\]

\[
(ii) \varepsilon \left(y _ {1}\right) + \dots + \varepsilon \left(y _ {i}\right) \geqslant 0 \quad \text { for } 1 \leqslant i \leqslant p - 1.
\]

证明 L - X 的每个元素都可以唯一地写成 cuv 的形式，其中 \(u, v \in L\) 。

（满足 (i) 与 (ii) 的元素 \(w\) 构成一个子广群 \(L'\) ，它含于 \(M\) 且包含 \(X\) 。若 \(w = y_1 \ldots y_p\) 属于 \(L'\) ，设 \(k\) 为使 \(\varepsilon(y_1) + \cdots + \varepsilon(y_k) = 0\) 的最小整数；证明 \(y_1 = c\) ，且元素 \(u = y_2 \ldots y_k\) 与 \(v = y_{k+1} \ldots y_p\) 属于 \(L'\) ；于是 \(w = cuv\) ，从而对 \(p\) 作归纳即得关系 \(w \in L\) 。然后证明，关系 \(u', v' \in L\) 与 \(w = cu'v'\) 蕴含 \(u' = u, v' = v\) 。）

\begin{enumerate}
\def\labelenumi{(\alph{enumi})}
\setcounter{enumi}{1}
\tightlist
\item
  证明 X 到 L 中的单射可延拓为 M(X) 到 L 上的同构。
\end{enumerate}

\begin{enumerate}
\def\labelenumi{\arabic{enumi}.}
\setcounter{enumi}{4}
\tightlist
\item
  设 X 为单元素集。若 n 为整数 \(\geqslant1\) ，设 \(u_{n}\) 表示 M(X) 中长度为 n 的元素个数。
\end{enumerate}

\begin{enumerate}
\def\labelenumi{(\alph{enumi})}
\tightlist
\item
  对于每一个集合Y，证明
\end{enumerate}

\[
\operatorname{Card} \left(\mathrm{M} _ {n} (\mathrm{Y})\right) = u _ {n}. \operatorname{Card} (\mathrm{Y}) ^ {n}.
\]

（利用习题 3。）

\begin{enumerate}
\def\labelenumi{(\alph{enumi})}
\setcounter{enumi}{1}
\tightlist
\item
  建立关系
\end{enumerate}

\[
u _ {n} = \sum_ {p = 1} ^ {n - 1} u _ {p} u _ {n - p} \quad \text { for } n \geqslant 2.
\]

\begin{enumerate}
\def\labelenumi{(\alph{enumi})}
\setcounter{enumi}{2}
\tightlist
\item
  *设 \(f(\mathbf{T})\) 是形式幂级数 \(\sum_{n=1}^{\infty} u_n \mathbf{T}^n\) 。证明
\end{enumerate}

\[
f (\mathbf {T}) = \mathbf {T} + f (\mathbf {T}) ^ {2}.
\]

由此推出公式

\[
f (\mathrm{T}) = \frac {1}{2} (1 - (1 - 4 \mathrm{T}) ^ {1 / 2}), \quad u _ {n} = \frac {2 ^ {n - 1}}{n !} 1. 3. 5 \dots (2 n - 3) \quad \text { for } n \geqslant 2. *
\]

利用《集合论》III，§ 5，习题 11 得到最后这一结果。

¶ 6. 设 X 为一个集合。

\begin{enumerate}
\def\labelenumi{(\alph{enumi})}
\item
  设 N 为 M(X) 的子广群，且 Y = N - (N.N)。单射 Y → N 可延拓为一个同态 u: M(Y) → N。证明 u 是同构。（换言之，自由广群的每个子广群都是自由的。）
\item
  若 \(x \in \mathbf{M}(\mathbf{X})\) ，设 \(\mathbf{M}_x\) 表示 \(\mathbf{M}(\mathbf{X})\) 中由 \(x\) 生成的子广群；由 (a)， \(\mathbf{M}_x\) 等同于在 \(x\) 上构造的自由广群。证明，若 \(x, y \in \mathbf{M}(\mathbf{X})\) ，则要么 \(\mathbf{M}_x \subset \mathbf{M}_y\) ，要么 \(\mathbf{M}_y \subset \mathbf{M}_x\) ，要么 \(\mathbf{M}_x \cap \mathbf{M}_y = \varnothing\) 。（若 \(\mathbf{M}_x \cap \mathbf{M}_y \neq \varnothing\) ，设 \(z\) 为 \(\mathbf{M}_x \cap \mathbf{M}_y\) 中长度最小的元素；证明 \(z = x\) 或 \(z = y\) 。）
\end{enumerate}

\begin{enumerate}
\def\labelenumi{\arabic{enumi}.}
\setcounter{enumi}{6}
\tightlist
\item
  设 \(\mathbf{M}_x\) 为在单元素集 \(\{x\}\) 上构造的自由广群。(a) 对所有 \(y \in \mathbf{M}_x\) ，证明存在唯一的自同态 \(f_y\) （\(\mathbf{M}_x\) 的），使得 \(f_y(x) = y\) ；它是从 \(\mathbf{M}_x\) 到子广群 \(\mathbf{M}_y\) （由 \(y\) 生成）上的同构（参见习题 6）。
\end{enumerate}

\begin{enumerate}
\def\labelenumi{(\alph{enumi})}
\setcounter{enumi}{1}
\item
  若 \(y, z \in \mathbf{M}_x\) ，记 \(y \circ z = f_y(z)\) 。证明合成律 \((y, z) \mapsto y \circ z\) 使 \(\mathbf{M}_x\) 成为一个以 \(x\) 为单位元的幺半群，它同构于幺半群 \(\operatorname{End}(\mathbf{M}_x)\) ——\(\mathbf{M}_x\) 的自同态所组成的幺半群。
\item
  设 \(\mathbf{N}\) 为 \(\mathbf{M}_x\) 的异于 \(\mathbf{M}_x\) 的单生成子广群组成的集合。一个元素 \(y\) （属于 \(\mathbf{M}_x\) ）称为本原的，如果 \(\mathbf{M}_y\) 是 \(\mathbf{N}\) 的一个极大元。证明 \(xx, x(x(xx))\) 是本原的，而 \(x, (xx)(xx)\) 不是。
\item
  设 \(\mathbf{P}\) 为 \(\mathbf{M}_x\) 的本原元组成的集合。证明，若 \(y, z \in \mathbf{P}, y \neq z\) ，则 \(\mathbf{M}_y \cap \mathbf{M}_z = \varnothing\) （利用习题 6）。
\item
  设 \(\operatorname{Mo}(\mathbf{P})\) 为在 \(\mathbf{P}\) 上构造的自由幺半群。若给 \(\mathbf{M}_x\) 赋予 (b) 中定义的幺半群结构，则单射 \(\mathbf{P} \to \mathbf{M}_x\) 可延拓为一个同态 \(p: \operatorname{Mo}(\mathbf{P}) \to \mathbf{M}_x\) 。证明 \(p\) 是一个同构。
\end{enumerate}

（若 \(z \in \mathbf{M}_x\) ，通过对 \(l(z)\) 作归纳证明 \(z \in \operatorname{Im}(p)\) 。另一方面，若 \(y_1, \ldots, y_n, z_1, \ldots, z_m \in \mathbb{P}\) 满足 \(p(y_1 \ldots y_n) = p(z_1 \ldots z_m)\) ，则 \(\mathbf{M}_{y_1} \cap \mathbf{M}_{z_1} \neq \varnothing\) ，从而 \(y_1 = z_1\) 且 \(f_{y_1}(p(y_2 \ldots y_n)) = f_{y_1}(p(z_2 \ldots z_m))\) ；由此得 \(p(y_2 \ldots y_n) = p(z_2 \ldots z_m)\) ，从而推出 \(y_1 \ldots y_n = z_1 \ldots z_m\) （对 \(\sup(n, m)\) 作归纳）。）

\begin{enumerate}
\def\labelenumi{\arabic{enumi}.}
\setcounter{enumi}{7}
\tightlist
\item
  设 X 为一个集合；对每个整数 \(q \geqslant 0\) ，设 \(\mathrm{Mo}(\mathbf{X})_{q}\) 为 \(\mathrm{Mo}(\mathbf{X})\) 中长度为 q 的元素组成的集合，且设 \(\mathrm{Mo}^{(q)}(\mathbf{X})\) 为诸 \(\mathrm{Mo}(\mathbf{X})_{nq}, n \in \mathbb{N}\) 的并。单射 \(\mathrm{Mo}(\mathbf{X})_{q} \to \mathrm{Mo}^{(q)}(\mathbf{X})\) 可延拓为一个同态
\end{enumerate}

\[
\operatorname{Mo} (\operatorname{Mo} (\mathrm{X}) _ {q}) \rightarrow \operatorname{Mo} ^ {(q)} (\mathrm{X});
\]

证明当 \(q \geqslant 1\) 时该同态是双射。

\begin{enumerate}
\def\labelenumi{\arabic{enumi}.}
\setcounter{enumi}{8}
\item
  设 \(x, y \in \mathbf{X}\) 且 \(x \neq y\) 。证明 \(\operatorname{Mo}(\mathbf{X})\) 中由 \(\{x, xy, yx\}\) 生成的子幺半群不同构于任何自由幺半群。
\item
  证明由表现
\end{enumerate}

\[
\langle x, y; x y ^ {2} = y ^ {3} x, y x ^ {2} = x ^ {3} y \rangle
\]

定义的群约化为单位元。

（第一个关系蕴含 \(x^{2}y^{8}x^{-2} = y^{18}\) 与 \(x^{3}y^{8}x^{-3} = y^{27}\) ；利用第二个关系可推出 \(y^{18} = y^{27}\) ，从而 \(y^{9} = e\) ；于是 \(y^{2}\) 共轭于 \(y^{3}\) 这一事实蕴含 \(y = e\) ，从而 \(x = e\) 。）

\begin{enumerate}
\def\labelenumi{\arabic{enumi}.}
\setcounter{enumi}{10}
\tightlist
\item
  由表现 \(\langle x,y;x^{2},y^{3},(xy)^{2}\rangle\) 定义的群同构于 \(S_{3}\) 。
\end{enumerate}

¶ 12. 由表现 \(\langle x, y; x^3, y^2, (xy)^3 \rangle\) 定义的群同构于 \(\mathfrak{A}_4\) 。（证明 \(y\) 的共轭元与单位元一起构成一个阶 \(\leqslant 4\) 、指数 \(\leqslant 3\) 的正规子群。）

\begin{enumerate}
\def\labelenumi{\arabic{enumi}.}
\setcounter{enumi}{12}
\tightlist
\item
  设 \(G\) 为一个群，\(X\) 为 \(G\) 的一个与 \(X^{-1}\) 不相交的子集。设 \(S = X \cup X^{-1}\) 。证明下列两个性质的等价性：
\end{enumerate}

\begin{enumerate}
\def\labelenumi{(\roman{enumi})}
\item
  X 是 G 中的自由族。
\item
  没有乘积 \(s_1 \ldots s_n\) （满足 \(n \geqslant 1\) , \(s_i \in S\) 且 \(s_i s_{i+1} \neq e\) 对 \(1 \leqslant i \leqslant n-1\) 成立）等于 \(e\) 。
\end{enumerate}

\begin{enumerate}
\def\labelenumi{\arabic{enumi}.}
\setcounter{enumi}{13}
\tightlist
\item
  群 G 称为自由群，如果它拥有一个基本族，从而同构于在某个集合 X 上构造的自由群 F(X)。
\end{enumerate}

\begin{enumerate}
\def\labelenumi{(\alph{enumi})}
\tightlist
\item
  证明，若 \((t_i)_{i\in I}\) 与 \((t_j')_{j\in J}\) 是两个这样的族，则 \(\operatorname{Card}(I) = \operatorname{Card}(J)\) 。（当 \(\operatorname{Card}(I)\) 无限时，证明 \(\operatorname{Card}(J)\) 也无限，且二者都等于 \(\operatorname{Card}(G)\) 。当 \(\operatorname{Card}(I)\) 是整数 \(d\) 时，证明 \(G\) 中指数为 2 的子群个数是 \(2^d - 1\) 。）
\end{enumerate}

G 的基本族的基数称为 G 的秩。

\begin{enumerate}
\def\labelenumi{(\alph{enumi})}
\setcounter{enumi}{1}
\item
  一个自由群的秩为 0（相应地，1）当且仅当它约化为 \(\{e\}\) （相应地，同构于 \(\mathbf{Z}\) ）。
\item
  设 G 为秩为 d 的自由群，且 \((x_{1},\ldots,x_{d})\) 为 G 的由 d 个元素组成的生成族。证明该族是基本的。（利用习题 34 与 § 5 习题 5。）
\item
  证明秩 \(\geqslant 2\) 的自由群包含一个具有给定秩 \(d\) 的自由子群，这里 \(d \leqslant \aleph_0\) 为任意基数。（利用习题 22。）
\end{enumerate}

¶ 15. 设 G 是一个具有表现 (t, r) 的群，其中 \(\mathbf{t} = (t_{i})_{i \in I}\) 与 \(\mathbf{r} = (r_{j})_{j \in J}\) ；令 \(\mathrm{F}((\mathrm{T}_{i})_{i \in I})\) 表示自由群 \(\mathrm{F}(I)\)，并令 \(\mathbf{T} = (\mathrm{T}_{i})\) 。另一方面，设 \(x = (x_{\alpha})_{\alpha \in A}\) 是 G 的一个生成族；记

\[
\mathrm{F} (\mathbf {A}) = \mathrm{F} ((\mathbf {X} _ {\alpha}) _ {\alpha \in A}),
\]

\(\mathbf{X} = (\mathbf{X}_{\alpha})\) ，而 \(\mathbf{R}_{\mathbf{x}}\) 表示 \(\mathbf{x}\) 的关系子集；它是 \(\mathrm{F(A)}\) 的一个正规子群。对所有 \(i \in \mathbf{I}\) ，设 \(\phi_i(\mathbf{X})\) 为 \(\mathrm{F(A)}\) 中使 \(t_i = \phi_i(\mathbf{x})\) 在 \(\mathbf{G}\) 中成立的元素；对所有 \(\alpha \in \mathbf{A}\) ，设 \(\psi_{\alpha}(\mathbf{T})\) 为 \(\mathrm{F(I)}\) 中使 \(x_{\alpha} = \psi_{\alpha}(\mathbf{t})\) 在 \(\mathbf{G}\) 中成立的元素。证明 \(\mathbf{R}_{\mathbf{x}}\) 是 \(\mathrm{F(A)}\) 中由元素 \(\mathbf{X}_{\alpha}^{-1} \psi_{\alpha}(\phi_i(\mathbf{X})_{i \in I}), \alpha \in \mathbf{A}\) 与 \(r_j(\phi_i(\mathbf{X})_{i \in I}) j \in J\) 生成的正规子群。（若 \(\mathbf{R}_{\mathbf{x}}'\) 表示 \(\mathrm{F(A)}\) 中由上述元素生成的正规子群，则 \(\mathbf{R}_{\mathbf{x}}' \subset \mathbf{R}_{\mathbf{x}}\) ；另一方面，证明同态 \(\mathrm{F(I)} \to \mathrm{F(A)}\) （由诸 \(\phi_i\) 所定义）在过渡到商时给出一个同态 \(\mathrm{G} \to \mathrm{F(A)} / \mathrm{R}_{\mathbf{x}}'\) ，它是典范同态 \(\mathrm{F(A)} / \mathrm{R}_{\mathbf{x}}' \to \mathrm{F(A)} / \mathrm{R}_{\mathbf{x}} = \mathrm{G}\) 的逆。）

\begin{enumerate}
\def\labelenumi{\arabic{enumi}.}
\setcounter{enumi}{15}
\tightlist
\item
  一个群称为有限生成的（相应地，有限表现的），如果它有一个表现 \(\langle (t_i)_{i\in I}, (r_j)_{j\in J}\rangle\) ，其中 \(I\) 是有限集（相应地，其中 \(I\) 与 \(J\) 都是有限集）。
\end{enumerate}

\begin{enumerate}
\def\labelenumi{(\alph{enumi})}
\item
  设 G 是一个有限表现的群，且 \(\mathbf{x} = (x_{\alpha})_{\alpha \in \mathbf{A}}\) 是 G 的一个有限生成族。证明存在族 x 的关系子的一个有限族 \(\mathbf{s} = (s_{k})_{k \in \mathbf{K}}\) ，使得 \((\mathbf{x}, \mathbf{s})\) 是 G 的一个表现。（利用前一习题。）
\item
  证明有限群是有限表现的。
\item
  给出一个有限表现的群的例子，它具有一个不是有限生成的子群。
\item
  若 \(\mathbf{G}\) 是一个有限表现的群，而 \(\mathbf{H}\) 是 \(\mathbf{G}\) 的一个正规子群，证明下列性质的等价性：
\item
  G/H 是有限表现的。
\end{enumerate}

\begin{enumerate}
\def\labelenumi{(\roman{enumi})}
\setcounter{enumi}{1}
\item
  存在 G 的有限子集 X，使得 H 是 G 中由 X 生成的正规子群。
\item
  若 \((\mathbf{H}_i)_{i\in \mathbf{I}}\) 是 \(\mathbf{G}\) 的一个右有向正规子群族，其并等于 \(\mathbf{H}\) ，则存在 \(i\in \mathbf{I}\) ，使得 \(\mathbf{H}_i = \mathbf{H}\) 。
\end{enumerate}

（利用 (a) 证明 (i) 蕴含 (ii)。）

\begin{enumerate}
\def\labelenumi{(\alph{enumi})}
\setcounter{enumi}{4}
\item
  若 G 是一个有限表现的群，证明 \(\mathbf{G}/\mathbf{C}^{r}(\mathbf{G})\) 对所有 \(r \geqslant 1\) 都是有限表现的。（利用 § 6 习题 15。）
\item
  设 \((\mathbf{G}_{\alpha})_{\alpha \in \mathbf{A}}\) 为一族群，且设 \(\mathbf{G} = \prod_{\alpha \in \mathbf{A}} \mathbf{G}_{\alpha}\) 为诸 \(\mathbf{G}_{\alpha}\) 的积。证明 \(\mathbf{G}\) 是有限生成的（相应地，有限表现的）当且仅当每个 \(\mathbf{G}_{\alpha}\) 都是有限生成的（相应地，有限表现的），且 \(\mathbf{G}_{\alpha} = \{e\}\) 对除有限个以外的 \(\alpha\) 成立。
\end{enumerate}

¶ 17. (a) 证明有限生成的幂零群的每一个子群都是有限生成的。（在交换情形中，通过对群的最小生成元个数进行归纳来论证；然后对群的幂零类进行归纳。）

给出一个有限生成的可解群的例子，它包含一个不是有限生成的子群。

\begin{enumerate}
\def\labelenumi{(\alph{enumi})}
\setcounter{enumi}{1}
\tightlist
\item
  证明每个有限生成的幂零群都是有限表现的。（将该群写成形式 \(\mathrm{F}(\mathbf{X}) / \mathbb{R}\) ，其中 \(\mathbf{X}\) 有限，并选取 \(r\) 使得 \(\mathbf{R} \supset \mathbf{C}^r (\mathrm{F}(\mathbf{X}))\) ；利用 (a) 证明 \(\mathbb{R} / \mathbb{C}^r (\mathrm{F}(\mathbf{X}))\) 是有限生成的；注意 \(\mathrm{F}(\mathbf{X}) / \mathbb{C}^r (\mathrm{F}(\mathbf{X}))\) 是有限表现的，参见习题 16。）
\end{enumerate}

\begin{enumerate}
\def\labelenumi{\arabic{enumi}.}
\setcounter{enumi}{17}
\tightlist
\item
  设 S 是群 G 的一个对称子集。G 的两个元素 \(g_{1}, g_{2}\) 称为 S-邻元，如果 \(g_{2}^{-1}g_{1}\) 属于 S。G 的元素的一个序列 \((g_{1}, \ldots, g_{n})\) 称为 S-链，如果 \(g_{i}\) 与 \(g_{i+1}\) 是 S-邻元（对 \(1 \leqslant i \leqslant n - 1\) ）；元素 \(g_{1}\) 与 \(g_{n}\) 分别称为链的起点与终点。
\end{enumerate}

\begin{enumerate}
\def\labelenumi{(\alph{enumi})}
\item
  设 Y 是 G 的一个子集，\(R_{Y}\{a, b\}\) 表示关系“在 Y 中存在一条从 a 开始而在 \(b''\) 结束的 S-链”。证明 \(R_{Y}\) 是 Y 上的一个等价关系（参见《集合论》II，§ 6，习题 10）；该关系下的等价类称为 Y 的连通 S-分量。若 Y 至多有一个连通分量，则称 Y 是 S-连通的。
\item
  证明 \(\mathbf{G}\) 是 S-连通的当且仅当 \(\mathbf{S}\) 生成 \(\mathbf{G}\) 。
\item
  假设没有元素 \(s \in S\) 满足关系 \(s^2 = e\) ；选取子集 \(X\) （含于 \(S\) ），使得 \(S\) 是 \(X\) 与 \(X^{-1}\) 的不交并。证明下列条件的等价性：
\item
  族 \(\mathbf{X}\) 是自由的。
\end{enumerate}

\begin{enumerate}
\def\labelenumi{(\roman{enumi})}
\setcounter{enumi}{1}
\tightlist
\item
  不存在 G 中的 S-链 \((g_{1},\ldots ,g_{n})\) ——由 \(n\geqslant 3\) 个互异元素组成——使得 \(g_{1}\) 与 \(g_{n}\) 是 S-邻元。
\end{enumerate}

（利用习题 13。）

¶ 19. 设 X 为一个集合，G = F(X) 为在 X 上构造的自由群，且 S = X ∪ X \(^{-1}\) 。

\begin{enumerate}
\def\labelenumi{(\alph{enumi})}
\item
  若 \(g \in G\) ，每个序列 \((s_1, \ldots, s_n)\) （由 \(S\) 的元素组成，满足 \(g = s_1 \ldots s_n\) 且 \(s_i s_{i+1} \neq e\) 对 \(1 \leqslant i < n\) 成立）称为 \(g\) 的既约分解。证明 \(G\) 的每个元素有且仅有一个既约分解；整数 \(n\) 称为 \(g\) 的长度（参见习题 26）。
\item
  设 Y 是 G 的包含 e 的子集。证明以下条件的等价性：
\end{enumerate}

\((b_{1})\) Y 是 S-连通的（参见习题 18）。

\((b_{2})\) 对所有 \(g \in \mathbf{Y}\) ，若 \((s_1, \ldots, s_n)\) 是 \(g\) 的既约分解，则 \(s_1 \ldots s_i \in \mathbf{Y}\) 对 \(1 \leqslant i \leqslant n\) 成立。

\begin{enumerate}
\def\labelenumi{(\alph{enumi})}
\setcounter{enumi}{2}
\tightlist
\item
  设 \(Y_{1}\) 与 \(Y_{2}\) 是 G 的两个非空的 S-连通子集，且 \(Y_{1} \cap Y_{2} = \varnothing\) 。证明至多存在一个有序对
\end{enumerate}

\[
\left(y _ {1}, y _ {2}\right) \in \mathbf {Y} _ {1} \times \mathbf {Y} _ {2}
\]

使得 \(y_{1}\) 与 \(y_{2}\) 是 S-邻元。这样的有序对存在当且仅当 \(\mathbf{Y}_1 \cup \mathbf{Y}_2\) 是 S-连通的；这时称 \(\mathbf{Y}_1\) 与 \(\mathbf{Y}_2\) 相邻。

\begin{enumerate}
\def\labelenumi{(\alph{enumi})}
\setcounter{enumi}{3}
\tightlist
\item
  设 \(Y_{1}, \ldots, Y_{n}\) 是 G 的一列互不相交的非空 S-连通子集。假设 \(Y_{i}\) 与 \(Y_{i+1}\) 相邻（对 \(1 \leqslant i < n\) ）。证明，若 \(n \geqslant 3\) ，则 \(Y_{1}\) 与 \(Y_{n}\) 不相邻。
\end{enumerate}

¶ 20. 我们保留前一习题的记号。设 H 是 G 的一个子群。

\begin{enumerate}
\def\labelenumi{(\alph{enumi})}
\item
  设 \(R_{H}\) 为 G 中满足 \(hT \cap T = \varnothing\) （若 \(h \in H, h \neq e\) ）的 S-连通子集 T 组成的集合。证明 \(R_{H}\) 按包含关系是归纳的。
\item
  设 Y 为 \(R_{H}\) 的一个极大元。证明 \(G = \bigcup_{h \in H} hY\) ，换言之，Y 是模 H 的右陪集的一个代表系。（若 \(G' = \bigcup hY\) 不同于 G，证明存在 \(x \in G'\) 与 \(y \in Y\) ，二者是 S-邻元，并由此推出 \(Y \cup \{x\}\) 属于 \(R_{H}\) 。）
\item
  设 \(S_{H}\) 为由使 Y 与 hY 相邻的元素 \(h \in H\) 组成的集合。它是 H 的一个对称子集。证明元素 \(h \in H\) 属于 \(S_{H}\) 当且仅当 \(h \neq e\) 且存在 \(y, y' \in Y, s \in S\) 使得 \(ys = hy'\) 。
\end{enumerate}

证明 \(S_{H}\) 生成 H。（若 \(H'\) 是由 \(S_{H}\)

生成的 H 的子群，证明 \(\bigcup_{h\in H'}hY\) 是 G 的一个连通 S-分量。）通过 § 4 习题 14 得到这一结果。

\begin{enumerate}
\def\labelenumi{(\alph{enumi})}
\setcounter{enumi}{3}
\tightlist
\item
  设 \(X_{H}\) 为 \(S_{H}\) 的一个子集，使得 \(S_{H}\) 是 \(X_{H}\) 与 \(X_{H}^{-1}\) 的不交并（证明这样的子集存在）。
\end{enumerate}

证明 \(X_{H}\) 是 H 的一个基本族，特别地，H 是自由群（“尼尔森–施赖尔定理”）。（将习题 18(c) 的判据应用于 H 及 \(S_{H}\) ，注意 H 的两个元素 h 与 \(h'\) 是 \(S_{H}\) -邻元当且仅当 hY 与 \(h'Y\) 是 G 的相邻子集。利用习题 19(d) 证明 \(S_{H}\) 满足习题 18(c) 中的条件 (ii)。）

\begin{enumerate}
\def\labelenumi{(\alph{enumi})}
\setcounter{enumi}{4}
\tightlist
\item
  设 \(d = (\mathbf{G}:\mathbf{H}) = \operatorname{Card}(\mathbf{Y})\) 。假设 \(d\) 是有限的。证明：
\end{enumerate}

\[
\begin{array}{l l} \operatorname{Card} (\mathbf {X} _ {\mathrm{H}}) = \operatorname{Card} (\mathbf {X}) = \operatorname{Card} (\mathbf {G}) & \text { if } \operatorname{Card} (\mathbf {X}) \text { is   infinite } \\ \operatorname{Card} (\mathbf {X} _ {\mathrm{H}}) = 1 + d (\operatorname{Card} (\mathbf {X}) - 1) & \text { if } \operatorname{Card} (\mathbf {X}) \text { is   finite }. \end{array}
\]

（设 T 是 G 的一个非空 S-连通有限子集，T'（相应地 T''）是其第一个元素（相应地两个元素）属于（相应地都属于）T 的 S-相邻元素的有序对组成的集合。Card(T') = 2Card(X)Card(T)，且应通过对 Card(T) 作归纳证明

\[
\operatorname{Card} \left(\mathrm{T} ^ {\prime \prime}\right) = 2 \operatorname{Card} (\mathrm{T}) - 2.
\]

将结果应用于 T = Y，注意 \(S_{H}\) 与 \(T' - T''\) 等势。

\begin{enumerate}
\def\labelenumi{\arabic{enumi}.}
\setcounter{enumi}{20}
\item
  设 H 是有限指数的群 G 的一个子群。证明 H 是有限表现的当且仅当 G 是有限表现的。（可假定 G 是有限生成的，参见 § 4 习题 14。将 G 写成形式 G = F/R，其中 F 是自由且有限生成的；H 在 F 中的原像 F' 是自由且有限生成的，参见习题 20，且 H = F'/R。若 R 作为 F 的正规子群由 \((r_{j})_{j\in J}\) 生成，证明它作为 F' 的正规子群由诸 \(yr_{j}y^{-1}\) , \(y\in Y\) , \(j\in J\) 生成，其中 Y 表示 F 中模 F' 的右陪集的一个代表系。）
\item
  \begin{enumerate}
  \def\labelenumii{(\alph{enumii})}
  \tightlist
  \item
    设 \((y_{n})_{n\in \mathbf{Z}}\) 是群 F 的一个基本族。设 \(\phi\) 为将 \(y_{n}\) 映为 \(y_{n + 1}\) 的 F 的自同构；该自同构定义了运算 \(\tau\) （\(\mathbf{Z}\) 在 F 上的运算）；设 \(\mathbf{E} = \mathbf{F}\times_{\tau}\mathbf{Z}\) 为相应的半直积。证明元素 \(x = (e,1)\) 与 \(y = (y_0,0)\) 构成 E 的一个基本族。（设 \(\mathrm{F}_{x,y}\) 为在 \(\{x,y\}\) 上构造的自由群，并设 \(f\) 为 \(\mathrm{F}_{x,y}\) 到 E 中的典范同态。证明存在同态 \(g\colon \mathbf{E}\to \mathbf{F}_{x,y}\) ，使得 \(g((0,1)) = x\) 且 \(g((y_n,0)) = x^n yx^{-n}\) ，并且 \(f\) 与 \(g\) 互为逆映射。）
  \end{enumerate}
\end{enumerate}

\begin{enumerate}
\def\labelenumi{(\alph{enumi})}
\setcounter{enumi}{1}
\item
  由此推出，\(F_{x,y}\) 的由 y 生成的正规子群是一个自由群，其基本族为 \((x^{n}yx^{-n})_{n\in\mathbf{Z}}\) 。将习题 20 应用于由诸 \(x^{n}, n\in Z\) 组成的代表系来得到这一结果。
\item
  将上述结论推广到任意秩的自由群。
\end{enumerate}

\begin{enumerate}
\def\labelenumi{\arabic{enumi}.}
\setcounter{enumi}{22}
\tightlist
\item
  设 \(\mathbf{F}\) 是基本族为 \((x, y)\) , \(x \neq y\) 的一个自由群。证明 F 的导群以换位子族 \((x^i, y^j)\) , \(i \in \mathbf{Z}\) , \(j \in \mathbf{Z}\) , \(i \neq 0\) , \(j \neq 0\) 为基本族。
\end{enumerate}

（利用习题 20 或习题 32。）

\begin{enumerate}
\def\labelenumi{\arabic{enumi}.}
\setcounter{enumi}{23}
\tightlist
\item
  设 G 是一个群，\(S = G - \{e\}\) 且 \(F = F(S)\) 。对所有 \(s \in S\) ，设 \(x_s\) 表示 F 中相应的元素。若 \(s, t \in S\) ，设 \(r_{s,t}\) 表示由下列方式定义的 F 中的元素：
\end{enumerate}

\[
r _ {s, t} = x _ {s} x _ {t} x _ {s t} ^ {- 1} \quad \text { if } s t \neq e \text { in } G
\]

\[
r _ {s, t} = x _ {s} x _ {t} \quad \text { if } s t = e \text { in } G.
\]

设 \(\phi\) 为将 \(x_{s}\) 映为 \(s\) 的、F 到 G 上的同态，并令 R 为 \(\phi\)

的核。证明族 \((r_{s,t})\) （\(s, t \in \mathbf{S} \times \mathbf{S}\)）是群 \(\mathbf{R}\) 的一个基本族。（将习题 20 应用于 \(\mathbf{G}\) 在 \(\mathbf{F}\) 中的由 \(e\) 与诸 \(x_{s'}\) 组成的代表系。）

\begin{enumerate}
\def\labelenumi{\arabic{enumi}.}
\setcounter{enumi}{24}
\tightlist
\item
  设 G 是一个群。证明以下性质的等价性：
\end{enumerate}

\begin{enumerate}
\def\labelenumi{(\alph{enumi})}
\item
  G 是一个自由群。
\item
  对于每个群 H 以及 G 被 H 的扩张 E，存在一个截面 \(G \rightarrow E\) .
\end{enumerate}

（为证明 (b) 蕴含 (a)，取 E 为一个自由群并使用习题 20。）

\begin{enumerate}
\def\labelenumi{\arabic{enumi}.}
\setcounter{enumi}{25}
\tightlist
\item
  设 \(\mathbf{X}\) 为一个集合，\(g\) 为自由群 \(\mathrm{F}(\mathbf{X})\) 的一个元素。设 \(g = \prod_{\alpha=1}^{n} x_{\alpha}^{m(\alpha)}\) 为 \(g\) 作为 \(\mathbf{X}\) 中元素的幂之乘积的典范分解（参见第 5 号，命题 7）。整数 \(l(g) = \sum_{\alpha} |m(\alpha)|\) 称为 \(g\) 的长度。
\end{enumerate}

\begin{enumerate}
\def\labelenumi{(\alph{enumi})}
\item
  \(g\) 称为循环既约的，如果或者 \(x_{1} \neq x_{n}\) ，或者 \(x_{1} = x_{n}\) 且 \(m(1)m(n) > 0\) 。证明 \(F(X)\) 的每个共轭类包含且仅包含一个循环既约元素；它是所讨论的类中长度最小的元素。
\item
  一个元素 \(x \in \mathrm{F}(\mathbf{X})\) 称为本原的，如果不存在整数 \(n \geqslant 2\) 和元素 \(g \in \mathrm{F}(\mathbf{X})\) 使得 \(x = g^n\) 。证明每个元素 \(z \neq e\) 在 \(\mathrm{F}(\mathbf{X})\) 中都可以唯一地写成形式 \(x^n\) ，其中 \(n \geqslant 1\) 且 \(x\) 是本原的。（将其归约为 \(z\) 是循环既约的情形，并注意到，若 \(z = x^n\) ，则元素 \(x\) 也是循环既约的。）
\end{enumerate}

证明 z 的中心化子与 x 的中心化子相同；它是由 x 生成的循环子群。

\begin{enumerate}
\def\labelenumi{\arabic{enumi}.}
\setcounter{enumi}{26}
\tightlist
\item
  设 \(G = \times_{A} G_{i}\) 为族 \((G_{i})_{i \in I}\) 的群沿公共子群 A 的融合和。对所有 \(i \in I\) ，选取一个子集 \(P_{i}\) ，含于 \(G_{i}\) 且满足第 3 号的条件 (A)。
\end{enumerate}

\begin{enumerate}
\def\labelenumi{(\alph{enumi})}
\tightlist
\item
  设 \(x \in G\) ，并设 \((a; i_1, \ldots, i_n; p_1, \ldots, p_n)\) 为 \(x\) 的一个既约分解；则
\end{enumerate}

\[
x = a \prod_ {\alpha = 1} ^ {n} p _ {\alpha}, \quad \text { with } \quad p _ {\alpha} \in \mathrm{P} _ {i _ {\alpha}} - \{e _ {i _ {\alpha}} \}, i _ {\alpha} \neq i _ {\alpha + 1}.
\]

证明，如果 \(i_1 \neq i_n\) ，则 \(x\) 的阶是无限的。证明，如果 \(i_1 = i_n\) ，则 \(x\) 共轭于一个其分解的长度 \(\leqslant n - 1\) 的元素。

\begin{enumerate}
\def\labelenumi{(\alph{enumi})}
\setcounter{enumi}{1}
\tightlist
\item
  由此推出，G 的每个有限阶元素都共轭于诸 \(G_{i}\) 之一的一个元素。特别地，如果诸 \(G_{i}\) 除 e 外没有有限阶元素，则 G 亦然。
\end{enumerate}

\begin{enumerate}
\def\labelenumi{\arabic{enumi}.}
\setcounter{enumi}{27}
\tightlist
\item
  我们保留前一练习的记号。
\end{enumerate}

\begin{enumerate}
\def\labelenumi{(\alph{enumi})}
\item
  设 \(\mathbf{N}\) 为 \(\mathbf{A}\) 的一个子群。假设对所有 \(i \in \mathbf{I}\) ，群 \(\mathbf{N}\) 在 \(\mathbf{G}_i\) 中正规。证明 \(\mathbf{N}\) 在 \(\mathbf{G} = \times_{\mathbf{A}} \mathbf{G}_i\) 中正规，且 \(\times_{\mathbf{A/N}} \mathbf{G}_i / \mathbf{N}\) 到 \(\mathbf{G}/\mathbf{N}\) 中的典范同态是一个同构。
\item
  对于所有 \(i \in \mathbf{I}\) ，设 \(\mathbf{H}_i\) 为 \(\mathbf{G}_i\) 的一个子群，并设 \(\mathbf{B}\) 为 \(\mathbf{A}\) 的一个子群。假设 \(\mathbf{H}_i \cap \mathbf{A} = \mathbf{B}\) 对所有 \(i\) 成立。证明 \(\star_B \mathbf{H}_i\) 到 \(\mathbf{G}\) 中的典范同态是单射；其像是 \(\mathbf{G}\) 的由诸 \(\mathbf{H}_i\) 生成的子群。
\end{enumerate}

¶ 29. 设 \(G = G_1 *_A G_2\) 为两个群 \(G_1\) 与 \(G_2\) 沿子群 A 的融合和。

\begin{enumerate}
\def\labelenumi{(\alph{enumi})}
\item
  证明若 \(G_{1}\) 与 \(G_{2}\) 都是有限生成的，则 G 是有限生成的。
\item
  假设 \(\mathbf{G}_1\) 与 \(\mathbf{G}_2\) 是有限表现的。证明下列两个性质的等价性：
\end{enumerate}

\begin{enumerate}
\def\labelenumi{(\arabic{enumi})}
\item
  G 是有限表现的。
\item
  A 是有限生成的。
\end{enumerate}

（首先证明 (1) 蕴含 (2)。另一方面，若 A 不是有限生成的，则存在 A 的子群的一个右有向族 \((\mathbf{A}_{i})_{i\in I}\) ，满足 \(\bigcup_{i\in I}\mathbf{A}_{i}=\mathbf{A}\) 且对所有 i 有 \(A_{i}\neq A\) 。设 \(H_{i}\) （相应地 H）为 \(G_{1}*G_{2}\) 到 \(G_{1}*_{A_{i}}G_{2}\) （相应地到 G）上的典范同态的核。证明 H 是诸 \(H_{i}\) 的并，且对所有 i 有 \(H_{i}\neq H\) ；然后应用习题 16(d)。）

\begin{enumerate}
\def\labelenumi{(\alph{enumi})}
\setcounter{enumi}{2}
\tightlist
\item
  由此给出一个有限生成但非有限表现的群的例子。（取 \(G_{1}\) 与 \(G_{2}\) 为秩 2 的自由群，A 为有限秩的自由群，参见习题 22。）
\end{enumerate}

\begin{enumerate}
\def\labelenumi{\arabic{enumi}.}
\setcounter{enumi}{29}
\tightlist
\item
  设 A 和 B 是两个群，G = A * B 为它们的自由积。
\end{enumerate}

\begin{enumerate}
\def\labelenumi{(\alph{enumi})}
\tightlist
\item
  设 \(\mathbf{Z}\) 为 \(\mathbf{A} \backslash \mathbf{G} / \mathbf{A}\) 中异于 \(\mathbf{A}\) 的一个元素。证明 \(\mathbf{Z}\) 包含且仅包含一个如下形式的元素 \(z\) ：
\end{enumerate}

\[
b _ {1} a _ {1} b _ {2} a _ {2} \dots b _ {n - 1} a _ {n - 1} b _ {n},
\]

其中 \(n \geqslant 1\) , \(a_i \in \mathbf{A} - \{e\}\) , \(b_i \in \mathbf{B} - \{e\}\) 。证明 \(\mathbf{Z}\) 的每个元素都可以唯一地写成形式 \(aza'\) ，其中 \(a, a' \in \mathbf{A}\) 。

\begin{enumerate}
\def\labelenumi{(\alph{enumi})}
\setcounter{enumi}{1}
\tightlist
\item
  设 \(x \in \mathbf{G} - \mathbf{A}\) 。设 \(f_x\) 为 \(\mathbf{A} * \mathbf{A}\) 到 \(\mathbf{G}\) 中的同态，它在 \(\mathbf{A} * \mathbf{A}\) 的第一个（相应地，第二个）因子上的限制是恒同映射（相应地，映射 \(a \mapsto xax^{-1}\) ）。证明 \(f_x\) 是单射。（将 \(x\) 写成上述形式 \(aza'\) ，从而化归为 \(x = z\) 的情形；利用 \(\mathbf{G}\) 中既约分解的唯一性。）
\end{enumerate}

特别地， \(\mathbf{A} \cap x\mathbf{A}x^{-1} = \{e\}\) ，且 \(\mathbf{A}\) 在 \(\mathbf{G}\) 中的正规化子是 \(\mathbf{A}\) 。

\begin{enumerate}
\def\labelenumi{(\alph{enumi})}
\setcounter{enumi}{2}
\tightlist
\item
  设 \(\mathbf{T}\) 为 \(\mathbf{A} \backslash \mathbf{G} / \mathbf{B}\) 的一个元素。证明 \(\mathbf{T}\) 包含且仅包含一个如下形式的元素 \(t\) ：
\end{enumerate}

\[
b _ {1} a _ {1} b _ {2} a _ {2} \dots b _ {n - 1} a _ {n - 1},
\]

其中 \(n \geqslant 1\) , \(a_i \in \mathbf{A} - \{e\}\) , \(b_i \in \mathbf{B} - \{e\}\) 。证明 \(\mathbf{T}\) 的每个元素都可以唯一地写成 atb 的形式，其中 \(a \in \mathbf{A}\) , \(b \in \mathbf{B}\) 。

\begin{enumerate}
\def\labelenumi{(\alph{enumi})}
\setcounter{enumi}{3}
\tightlist
\item
  设 \(x \in G\) ，并设 \(h_x\) 为 \(G\) 的唯一自同态，使得 \(h_x(a) = a\) 当 \(a \in A\) ，且 \(h_x(b) = x b x^{-1}\) 当 \(b \in B\) 。证明 \(h_x\) 是单射。（用与 (b) 中相同的方法，利用 (c) 中的分解 \(x = a t b\) 。）特别地， \(A \cap x B x^{-1} = \{e\}\) 。
\end{enumerate}

给出一个 \(h_{x}\) 不是满射的例子。

\begin{enumerate}
\def\labelenumi{\arabic{enumi}.}
\setcounter{enumi}{30}
\tightlist
\item
  设 G 为某个群的族 \((G_{i})_{i\in I}\) 的自由积。设 S 为 G 的与诸 \(G_{i}\) 之一共轭的子群组成的集合，并设 \(\Theta\) 为 S 中所有元素的并。
\end{enumerate}

\begin{enumerate}
\def\labelenumi{(\alph{enumi})}
\item
  设 \(\mathbf{H}, \mathbf{H}' \in \mathbf{S}\) ，其中 \(\mathbf{H} \neq \mathbf{H}'\) ，并设 \(x \in \mathbf{H} - \{e\}\) , \(x' \in \mathbf{H}' - \{e\}\) 。证明 \(xx' \in \mathbf{G} - \Theta\) 。（化归到 \(\mathbf{H}\) 是诸 \(\mathbf{G}_i\) 之一的情形；将 \(\mathbf{H}'\) 写成形式 \(y\mathrm{G}_j y^{-1}\) ，并像前一习题那样进行。）
\item
  设 \(\mathbf{C}\) 为 \(\mathbf{G}\) 的一个含于 \(\Theta\) 的子群。证明存在 \(\mathbf{H} \in \mathbf{S}\) ，使得 \(\mathbf{C} \subset \mathbf{H}\) 。
\item
  设 C 为 G 的一个其所有元素均为有限阶的子群。证明 C 含于 \(\Theta\) （参见习题 27）；由此推出 C 含于诸 \(G_{i}\) 之一的一个共轭之中。
\end{enumerate}

¶ 32. 设 A 和 B 是两个群，G = A * B 是它们的自由积。设 X 为形如 \((a, b)\) 、其中 \(a \in A - \{e\}\) 且 \(b \in B - \{e\}\) 的换位子组成的集合，并设 \(S = X \cup X^{-1}\) 。

\begin{enumerate}
\def\labelenumi{(\alph{enumi})}
\item
  证明 \(\mathbf{X} \cap \mathbf{X}^{-1} = \varnothing\) ，并且，若 \(x\) 是 \(\mathbf{X}\) 的一个元素，则存在唯一的有序对 \(a, b\) ，属于 \((\mathbf{A} - \{e\}) \times (\mathbf{B} - \{e\})\) ，使得 \((a, b) = x\) 。
\item
  设 \(n \geqslant 1\) ，并设 \(s_1, \ldots, s_n\) 为 \(S\) 中元素的一个序列，使得 \(s_i s_{i+1} \neq e\) 对 \(1 \leqslant i < n\) 成立；设 \(a_n \in A - \{e\}, b_n \in B - \{e\}, \varepsilon(n) = \pm 1\) 满足 \(s_n = (a_n, b_n)^{\varepsilon(n)}\) 。设 \(g = s_1 \ldots s_n\) 。通过对 \(n\) 作归纳证明： \(g\) 的既约分解的长度 \(\geqslant n + 3\) ，并且或者以 \(a_nb_n\) 结束（若 \(\varepsilon(n) = 1\)），或者以 \(b_na_n\) 结束（若 \(\varepsilon(n) = -1\)）。
\item
  设 R 为典范投影 \(A*B \rightarrow A \times B\) 的核，该投影（第 3 号）对应于典范单射 \(A \rightarrow A \times B\) 与 \(B \rightarrow A \times B\) 。证明 R 是以 X 为基本族的自由群。（利用 (b) 证明 X 是自由的；若 \(R_{X}\) 为 G 中由 X 生成的子群，证明 \(R_{X}\) 是正规的；由此推出它与 R 重合。）
\end{enumerate}

\begin{enumerate}
\def\labelenumi{\arabic{enumi}.}
\setcounter{enumi}{32}
\tightlist
\item
  设 \(E = G *_{A} G'\) 为两个群沿公共子群 A 的融合和。设 P（相应地 \(P'\) ）为 G（相应地 \(G'\) ）的满足第 3 号的条件 (A) 的子集。
\end{enumerate}

设 \(n\) 为一个整数 \(\geqslant 1\) 。设 \(L_n\) 表示 \(x \in E\) 中既约分解的长度 \(\leqslant 2n - 1\) 者组成的集合。设 \(\mathbf{M}_n\) 表示 \(x \in \mathbf{E}\) 中形如 \(ap_1p_1'p_2p_2'\ldots p_np_n'\) 的元素组成的集合，其中 \(a \in \mathbf{A}, p_i \in \mathbf{P} - \{e\}, p_i' \in \mathbf{P}' - \{e\}\) ；类似地，令 \(\mathbf{M}_n'\) 表示形如 \(ap_1'p_1p_2'p_2\ldots p_np_n'\) 的元素组成的集合，其中 \(a, p_i, p_i'\) 满足相同的条件。

\begin{enumerate}
\def\labelenumi{(\alph{enumi})}
\item
  证明 \(L_{n}\) , \(M_{n}\) 与 \(M_{n}^{\prime}\) 两两不相交，且它们的并是 E 中既约分解的长度 \(\leqslant2n\) 的元素组成的集合。
\item
  构造一个双射 \(\varepsilon\) ——\(\mathbf{M}_n\) 到 \(\mathbf{M}_n'\) 上的双射——使得 \(\varepsilon(ax) = a\varepsilon(x)\) （若 \(a \in \mathbf{A}\) , \(x \in \mathbf{M}_n\)）。
\item
  设 \(\mathbf{X}_n = \mathbf{L}_n \cup \mathbf{M}_n\) 与 \(\mathbf{X}_n' = \mathbf{L}_n \cup \mathbf{M}_n'\) ；证明 \(\mathbf{G} \cdot \mathbf{X}_n = \mathbf{X}_n\) 且 \(\mathbf{G}' \cdot \mathbf{X}_n' = \mathbf{X}_n'\) 。
\item
  \(\varepsilon\) 扩充为 \(\mathbf{X}_n\) 到 \(\mathbf{X}_n'\) 上的一个双射，扩充方式为令 \(\varepsilon(x) = x\) （若 \(x \in \mathbf{L}_n\)）。证明 E 可以如下作用在 \(\mathbf{X}_n\) 上： \(g(x) = gx\) （若 \(g \in G\) , \(x \in \mathbf{X}_n\)），且 \(g'(x) = \varepsilon^{-1}(g' \varepsilon (x))\) （若 \(g' \in G'\) , \(x \in \mathbf{X}_n\)）。证明，若 \(g \in \mathbf{X}_n\) ，则 \(g(e) = g\) 。
\item
  设 \(\tau_{n}\colon\mathbf{E}\to\mathfrak{S}_{\mathbf{X}_{n}}\) 为由上述 \(\mathbf{E}\) 在 \(\mathbf{X}_{n}\) 上的作用所定义的同态，并令 \(\mathbf{R}_{n}\) 为 \(\tau_{n}\) 的核。证明 \(\mathbf{R}_{n}\cap\mathbf{X}_{n}=\{e\}\) ；由此推出诸 \(\mathbf{R}_{n}\) 的交化为 \(e\) 。
\item
  当 G 与 G' 有限时，证明 E 是剩余有限的（参见 § 5 习题 5）。（注意此时诸 R \(_{n}\) 都是有限指数的子群。）
\end{enumerate}

¶ 34. (a) 设 \((\mathrm{H}_{i})\) 为群 G 的一个右有向正规子群族。假设 \(\bigcap_{i} \mathrm{H}_{i} = \{e\}\) 。证明：对每个群 \(\mathbf{G}'\) ，典范同态 \(\mathbf{G}' * \mathbf{G} \to \mathbf{G}' * (\mathbf{G}/\mathbf{H}_{i})\) 的核的交等于 \(\{e\}\) 。

\begin{enumerate}
\def\labelenumi{(\alph{enumi})}
\setcounter{enumi}{1}
\item
  设 G 为某个群的族 \((G_{\alpha})\) 的自由积。证明，如果所有 \(G_{\alpha}\) 都是剩余有限的（参见 § 5 习题 5），那么 G 也是。（首先化归为族有限的情形，然后化归为族有两个元素的情形；由 (a)，可设诸 \(G_{\alpha}\) 有限，再应用前一习题。）
\item
  由此推出，每个自由群都是剩余有限的。
\end{enumerate}

\begin{enumerate}
\def\labelenumi{\arabic{enumi}.}
\setcounter{enumi}{34}
\item
  *设 G（相应地，G'）为形如 \(a/b\) 、其中 \(a \in \mathbf{Z}\) , \(b \in \mathbf{Z}\) 且 \(b\) 不被 2（相应地，不被 3）整除的分数组成的加法群。设 E = G * \(_{\mathbf{Z}}\) G' 为 G 与 G' 沿其公共子群 Z 的融合和。证明 G 与 G' 是剩余有限的，但 E 不是。（证明 E 的有限指数子群只有 E 本身。）*
\item
  设 \((\mathbf{G}_1, \ldots, \mathbf{G}_n)\) 与 \((\mathbf{A}_1, \ldots, \mathbf{A}_{n-1})\) 为两个群列；对每个 \(i\) ，设 \(\phi_i: \mathbf{A}_i \to \mathbf{G}_i\) 与 \(\psi_i: \mathbf{A}_i \to \mathbf{G}_{i+1}\) 为两个单同态；通过这些同态，\(\mathbf{A}_i\) 既可与 \(\mathbf{G}_i\) 的一个子群等同，也可与 \(\mathbf{G}_{i+1}\) 的一个子群等同。
\end{enumerate}

\begin{enumerate}
\def\labelenumi{(\alph{enumi})}
\item
  设 \(H_{1}=G_{1}, H_{2}=H_{1} *_{A_{1}} G_{2}, \ldots, H_{n}=H_{n-1} *_{A_{n-1}} G_{n}\) 。群 \(H_{n}\) 称为群 \((G_{i})\) 沿子群 \((A_{i})\) 的融合和；记为 \(G_{1} *_{A_{1}} G_{2} *_{A_{2}} \cdots *_{A_{n-1}} G_{n}\) 。若 T 为任意群，在 \(H_{n}\) 到 T 中的同态与同态族 \((t_{1},\ldots ,t_{n})\) （其中 \(t_i\colon \mathbf{G}_i\to \mathbf{T}\)）之间定义一个双射，使得 \(t_i\circ \phi_i = t_{i + 1}\circ \psi_i\) 对 \(1\leqslant i\leqslant n - 1\) 成立。
\item
  将每个 \(G_{i}\) 与 \(H_{n}\) 中相应的子群等同。证明 \(G_{i} \cap G_{i+1} = A_{i}\) ，并且 \(G_{i} \cap G_{i+2} = A_{i} \cap A_{i+1}\) （后一交集取在 \(G_{i+1}\) 中）。
\end{enumerate}

¶ 37. 设 G 为一个群，S 为 G 的有限对称子集且生成 G。令 \(\mathfrak{T}\) 表示 G 的有限子集全体，按包含关系排序；它是一个有向集。

\begin{enumerate}
\def\labelenumi{(\alph{enumi})}
\item
  设 \(\mathbf{T} \in \mathfrak{T}\) ，并设 \(\mathbf{B}_{\mathbf{T}}\) 为 \(\mathbf{G} - \mathbf{T}\) 的连通 S-分量组成的集合（参见习题 18）；设 \(\mathbf{B}_{\mathbf{T}}^{\infty} = \mathbf{B}_{\mathbf{T}} - (\mathfrak{T} \cap \mathbf{B}_{\mathbf{T}})\) 。证明 \(\mathbf{B}_{\mathbf{T}}\) 是有限的（注意，若 \(\mathbf{T} \neq \varnothing\) 且 \(\mathbf{Z} \in \mathbf{B}_{\mathbf{T}}\) ，则存在 \(t \in \mathbf{T}\) 与 \(z \in \mathbf{Z}\) ，它们是 S-相邻的）。
\item
  设 \(\mathbf{T}, \mathbf{T}' \in \mathfrak{T}\) 满足 \(\mathbf{T} \subset \mathbf{T}'\) 。若 \(Z' \in \mathbf{B}_{\mathbf{T}'}\) ，证明存在唯一的元素 \(Z \in \mathbf{B}_{\mathbf{T}}\) 使得 \(Z' \subset Z\) ；映射 \(f_{\mathbf{T}\mathbf{T}'}: \mathbf{B}_{\mathbf{T}'} \to \mathbf{B}_{\mathbf{T}}\) （满足 \(f_{\mathbf{T}\mathbf{T}'}(\mathbf{Z}') = \mathbf{Z}\)）把 \(\mathbf{B}_{\mathbf{T}}^{\infty}\) 映到 \(\mathbf{B}_{\mathbf{T}}^{\infty}\) 上。诸 \(\mathbf{B}_{\mathbf{T}}\) （其中 \(\mathbf{T}\) 遍历 \(\mathfrak{T}\)）的逆向极限记作 \(\mathbf{B}\) 。证明 \(B = \varprojlim B_{\mathbf{T}}^{\infty}\) ，且 \(B\) 在 \(B_{\mathbf{T}}\) 中的像是 \(B_{\mathbf{T}}^{\infty}\) 。集合 \(B\) 称为 \(G\) 的端集；它非空当且仅当 \(G\) 是无限的。
\item
  设 \(\mathbf{T}\) 为 \(\mathbf{G}\) 的一个非空 S-连通有限子集，并设 \(\mathbf{Z} \in \mathbf{B}_{\mathbf{T}}^{\infty}\) 。证明存在 \(g \in \mathbf{G}\) ，使得 \(g\mathbf{T} \cap \mathbf{T} = \varnothing\) 且 \(g\mathbf{T} \cap \mathbf{Z} \neq \varnothing\) ，而此时 \(g\mathbf{T} \subset \mathbf{Z}\) 。证明存在 \(\mathbf{Z}_1 \in \mathbf{B}_{\mathbf{T}}\) ，使得 \(g\mathbf{Z}_1\) 包含 \(\mathbf{T}\) 以及所有的 \(\mathbf{Z}' \in \mathbf{B}_{\mathbf{T}}\) （满足 \(\mathbf{Z}' \neq \mathbf{Z}\) ）（注意 \(\mathbf{T} \cup \bigcup_{\mathbf{Z}' \neq \mathbf{Z}} \mathbf{Z}'\) 是 S-连通的且不与 \(g\mathbf{T}\) 相交）；由此推出，若 \(\mathbf{Z}' \in \mathbf{B}_{\mathbf{T}}\) 不同于 \(\mathbf{Z}_1\) ，则 \(g\mathbf{Z}' \subset \mathbf{Z}\) 。证明，记 \(\mathbf{T}' = \mathbf{T} \cup g\mathbf{T}\) ，则 \(\mathbf{Z}\) 在 \(\mathbf{B}_{\mathbf{T}}^{\infty}\) 中的原像至少由 \(n - 1\) 个元素组成，其中 \(n = \operatorname{Card}(\mathbf{B}_{\mathbf{T}}^{\infty})\) 。
\item
  设 \(S'\) 为 \(G\) 的一个有限对称子集，它生成 \(G\) ，并设 \(B'\) 为 \(G\) 的、相对于 \(S'\) 的端集。定义 \(B\) 到 \(B'\) 上的一个双射（化归为 \(S \subset S'\) 的情形）。于是 \(B\) 的基数与 \(S\) 的选取无关；它称为 \(G\) 的端数。
\item
  利用 (c) 证明，G 的端数等于 0、1、2 或 \(2^{\aleph_0}\) 。
\item
  证明 \(\mathbf{Z}\) 的端数是 2，而 \(\mathbf{Z}^n\) ( \(n \geqslant 2\) ) 的端数是 1。
\end{enumerate}

¶ 38. 设 X 为有限集，F(X) 为在 X 上构造的自由群，且 S = X ∪ X \(^{-1}\) 。

\begin{enumerate}
\def\labelenumi{(\alph{enumi})}
\item
  设 \(\mathbf{S}_n\) 为乘积 \(s_1 \ldots s_m\) , \(s_i \in \mathbf{S}\) , \(m \leqslant n\) 组成的集合。证明 \(\mathrm{F}(\mathbf{X}) - \mathbf{S}_n\) 的每个连通 S-分量（参见习题 18）包含且仅包含一个形如 \(s_1 \ldots s_{n+1}\) 的元素，其中 \(s_i \in \mathbf{S}\) 且 \(s_i s_{i+1} \neq e\) 对 \(1 \leqslant i \leqslant n\) 成立。
\item
  由 (a) 推出集合 B（即 F(X) 的端集，参见习题 37）到无穷序列 \((s_{1},\ldots,s_{n},\ldots)\) （其元素属于 S 且满足 \(s_{i}s_{i+1}\neq e\) 对所有 i 成立）组成的集合上的一个双射。特别地，F(X) 的端数等于 \(2^{\aleph_0}\) （若 \(\text{Card}(X)\geqslant2\)）。
\end{enumerate}

\begin{enumerate}
\def\labelenumi{\arabic{enumi}.}
\setcounter{enumi}{38}
\tightlist
\item
  *设 G 是一个有限生成群，F 是 G 的包含 e 且生成 G 的对称有限子集的集合。
\end{enumerate}

\begin{enumerate}
\def\labelenumi{(\alph{enumi})}
\tightlist
\item
  若 \(\mathbf{X} \in \mathfrak{F}\) 且 \(n \in \mathbf{N}\) ，设 \(d_n(\mathbf{X})\) 表示 \(\mathbf{G}\) 中形如 \(x_1 \ldots x_n\) （其中 \(x_i \in \mathbf{X}\)）的元素的个数。若 \(\mathbf{X}, \mathbf{Y} \in \mathfrak{F}\) ，证明存在一个整数 \(a \geqslant 1\) ，使得
\end{enumerate}

\[
d _ {n} (\mathbf {X}) \leqslant d _ {a n} (\mathbf {Y}) \quad \text { for   all } n \in \mathbf {N}.
\]

由此推出 \(\limsup_{n\to \infty}\frac{\log(d_n(\mathbf{X}))}{\log(n)}\) 不依赖于 \(\mathbf{X}\) 在 \(\mathfrak{F}\) 中的选取；这一极限（有限或无限）记作 \(e(\mathbf{G})\) 。

\begin{enumerate}
\def\labelenumi{(\alph{enumi})}
\setcounter{enumi}{1}
\item
  若 \(\mathbf{X} \in \mathfrak{F}\) 且 \(\mathbf{G}\) 是无限的，证明 \(d_n(\mathbf{X}) < d_{n+1}(\mathbf{X})\) 对所有 \(n\) 成立。由此推出 \(e(\mathbf{G})\) \(\geqslant 1\) 。
\item
  设 H 为一个有限生成群。证明 \(e(\mathbf{H}) \leqslant e(\mathbf{G})\) ，如果 H 同构于 G 的一个子群（相应地，一个商群）。若 E 是 G 被 H 的扩张，证明 \(e(\mathbf{E}) \geqslant e(\mathbf{H}) + e(\mathbf{G})\) ，且当 \(E = G \times H\) 时等号成立。
\item
  证明 \(e(\mathbf{Z}^n) = n\) 。
\item
  设 \(\mathbf{X} \in \mathfrak{F}\) 。考虑下列性质：
\item
  存在一个实数 \(a > 1\) ，使得 \(d_n(\mathbf{X}) \geqslant a^n\) 对所有充分大的 \(n\) 成立。
\end{enumerate}

证明：若性质 (i) 对 \(\mathfrak{F}\) 的某个元素成立，则它对 \(\mathfrak{F}\) 的每个元素都成立。此时称 \(\mathbf{G}\) 为指数型的；这蕴含 \(e(\mathbf{G}) = +\infty\) 。

\begin{enumerate}
\def\labelenumi{(\alph{enumi})}
\setcounter{enumi}{5}
\item
  设 \(G_{1}\) 与 \(G_{2}\) 为包含同一子群 A 的两个群，令 \(H = G_{1} *_{A} G_{2}\) 为相应的融合和。假设 \(G_{1}\) 与 \(G_{2}\) 是有限生成的（从而 H 也是），且对 i = 1, 2 有 \(G_{i} \neq A\) 。证明：若指标 ( \(G_{1}:A\) ), ( \(G_{2}:A\) ) 中至少有一个 \(\geqslant 3\) ，则 H 是指数型的。
\item
  证明秩 \(\geqslant 2\) 的自由群是指数型的。*
\end{enumerate}

\begin{enumerate}
\def\labelenumi{\arabic{enumi}.}
\setcounter{enumi}{39}
\tightlist
\item
  *设 \(n\) 为一个整数 \(\geqslant 2\) ，\(T_n\) 为对角项都等于 1 的 \(n\) 阶上三角矩阵（在 \(Z\) 上）组成的群。设 \(e(T_n)\) 为习题 39 中定义的群 \(T_n\) 的不变量。证明
\end{enumerate}

\[
e (\mathrm{T} _ {n}) \leqslant \sum_ {i = 1} ^ {n} i (n - i). *
\]

